%% 60-conclusion.tex — Conclusion

\section{Conclusion}
\label{sec:conclusion}

The semantic diversity of EDM metadata across institutional sectors cannot be
resolved by schema-level alignment alone.
Corpus profiling --- analysing the actual field distributions of a deployed
metadata corpus before constructing any alignment rule --- is a precondition
for reliable automated mapping.
The three-layer dispatch architecture presented here operationalises this
principle: property alignment via a data-derived DC\,$\to$\,RDA table
(Layer~1), archival hierarchy typing via a DDB hierarchy code lookup
(Layer~2), and document-type dispatch via a dc:type\,$\times$\,sector\,$\times$\,mediatype
table (Layer~3).

Applied to 115{,}432 DDB records spanning all seven institutional sectors,
the pipeline produces 44.4 million RDF triples aligned to \textsc{mocho}'s
WEMI-level hub ontology, with domain-specific types from RDA, RiC-O, VRA~Core,
and Music Ontology.
Known alignment gaps are documented explicitly in the output tables rather
than suppressed, enabling downstream users to reason about query completeness.

\paragraph{Future work.}
Three directions are planned.
First, the pipeline will be applied to the full 65M-record DDB corpus,
requiring a profiling pass to discover any new (entity type, JSON key) pairs
not seen in the pilot.
Second, Layer~1 will be extended to cover EDM-structural and agent authority
properties currently out of scope, using a PROV-O mapping for provenance
properties and direct RDA Person properties for agent attributes.
Third, NER-based enrichment over contributor strings and title fields will
augment Layer~1 property alignment with entity-level role disambiguation,
supporting finer-grained RDA sub-property assignment
(e.g., \texttt{rdaa:P50450} for editors, \texttt{rdaw:P10053} for composers).
